\subsection{Relatifizinbarkeit und Bogenlänge}

Ziel dieses Abschnitts ist es, die Länge einer Kurve $\gamma: [a,b] \to \mathbb{R}^n$ zu definieren. Dazu approximieren wir die Länge über Zerlegungen des Intervalls.

Sei 


\[
Z = \{t_0, t_1, \dots, t_n\} \quad \text{mit} \quad a = t_0 < t_1 < \cdots < t_n = b
\]


eine Zerlegung von $[a,b]$. Für diese Zerlegung definieren wir die approximative Bogenlänge durch


\[
L(Z,\gamma) = \sum_{i=1}^n \|\gamma(t_i)-\gamma(t_{i-1})\|.
\]



\begin{definition}{}
Eine Kurve $\gamma: [a,b] \to \mathbb{R}^n$ heißt \emph{rektifizierbar} (bzw. relativizierbar) mit Bogenlänge $L(\gamma)$, falls es eine Zahl $L(\gamma)\in [0,+\infty)$ gibt, so dass für alle $\epsilon>0$ ein $\delta>0$ existiert, sodass für jede Zerlegung $Z$ von $[a,b]$ mit 


\[
\|Z\| := \max_{1\le i\le n}(t_i-t_{i-1}) < \delta
\]


folgt


\[
\bigl|L(Z,\gamma)-L(\gamma)\bigr|<\epsilon.
\]


\end{definition}

\begin{bemerkungbox}
Wesentlich ist, dass sich $L(Z,\gamma)$ durch Verfeinern einer Zerlegung (also durch Hinzufügen weiterer Zwischenwerte) nicht verringert – es gilt:


\[
L(\gamma)=\sup\{\,L(Z,\gamma):Z \text{ ist eine Zerlegung von } [a,b]\,\}.
\]


Dies entspricht den Prinzipien von Ober- und Untersummen beim Riemann-Integral.
\end{bemerkungbox}

\begin{satz}{}
Sei $\gamma: [a,b] \to \mathbb{R}^n$ stetig differenzierbar. Dann ist $\gamma$ rektifizierbar und es gilt


\[
L(\gamma)=\int_a^b \|\gamma'(t)\|\,dt.
\]


\end{satz}

\begin{beweisbox}
Die Idee des Beweises ist, die Summe


\[
\sum_{i=1}^n \|\gamma(t_i)-\gamma(t_{i-1})\|
\]


über eine Zerlegung zu bestimmen und zu zeigen, dass im Limes (beim Verfeinern der Zerlegung) diese Summe gegen


\[
\int_a^b \|\gamma'(t)\|\,dt
\]


konvergiert. Zur Herleitung werden folgende Lemmata verwendet.
\end{beweisbox}

\begin{lemma}
Sei $\gamma: [a,b] \to \mathbb{R}^n$ stetig differenzierbar und seien $t_0,t\in[a,b]$ mit $t_0<t$. Dann existiert ein $c\in (t_0,t)$, so dass


\[
\|\gamma(t)-\gamma(t_0)\|=\|\gamma'(c)\|(t-t_0).
\]


\end{lemma}

\begin{beweisbox}
Der Beweis erfolgt, indem man für jede Komponentenfunktion den Mittelwertsatz anwendet und anschließend die Norm betrachtet. (Details entnehmen Sie bitte der entsprechenden Literatur.)
\end{beweisbox}

\begin{lemma}
Sei $f: [a,b] \to \mathbb{R}$ stetig. Dann existiert ein $c\in [a,b]$, so dass


\[
\int_a^b f(x)\,dx = f(c)(b-a).
\]


\end{lemma}

\begin{beweisbox}
Der Beweis dieses klassischen Satzes folgt direkt aus den Eigenschaften der Riemann-Integration.
\end{beweisbox}

\begin{beispielbox}
\begin{enumerate}
  \item \textbf{Helix:} Betrachte die Helix
  

\[
  \gamma(t)=\bigl( r\cos t,\, r\sin t,\, c\,t \bigr),\quad t\in[0,2\pi],
  \]


  mit $r>0$ und $c\neq0$. Dann ist
  

\[
  \gamma'(t)=\bigl( -r\sin t,\, r\cos t,\, c \bigr)
  \]


  und die Bogenlänge berechnet sich zu
  

\[
  L(\gamma)=\int_0^{2\pi}\sqrt{r^2+c^2}\,dt=2\pi\sqrt{r^2+c^2}.
  \]


  
  \item \textbf{Wichtiger Hinweis:} Die Bogenlänge misst die zurückgelegte Wegstrecke der Kurve und entspricht nicht zwingend der "geraden" geometrischen Entfernung zwischen Start- und Endpunkt.
\end{enumerate}
\end{beispielbox}

\begin{bemerkungbox}
\begin{enumerate}
  \item Die Bogenlänge einer Kurve ist unabhängig von der Parametrisierung, d.h. reparametrisiert man eine Kurve, ändert sich ihre Länge nicht.
  \item Insbesondere können reguläre Kurven so reparametrisiert werden, dass sie nach Bogenlänge parametrisiert sind (sogenannte Einheitsgeschwindigkeitsparametrisierung).
\end{enumerate}
\end{bemerkungbox}
